% =============================================================
% CHAPTER 7 — CONCLUSIONS
% Purpose: close the loop. Remind the reader what you set out to do, what you
% found, why it matters, and what comes next. No NEW results or figures here.
% Length: typically 2-4 pages.
% =============================================================
\chapter{Conclusions}
\label{chap:conclusion}

\section{Summary}
\guide{Recap the thesis in one tight paragraph: the problem, your approach, and
the headline findings. This should echo the introduction and abstract -- the
reader should feel the arc has closed.}

\section{Answers to the Research Questions}
\guide{Revisit each research question from \Cref{chap:introduction} and answer it
in one or two sentences each, pointing to the evidence in \Cref{chap:results}.
This is the most direct way to show the thesis did what it promised.}

\section{Contributions Revisited}
\guide{Briefly restate the contributions now that they are backed by results --
what can the community take away and build on?}

\section{Future Work}
\guide{Suggest concrete, credible next steps that follow from your limitations:
e.g. extending to video, fine-tuning MLLMs on detection, larger/newer model
coverage, combining MLLMs with dedicated detectors, or studying robustness to
adversarial manipulation. Make these specific enough that someone could pick one
up.}
